\documentclass[10pt]{article}
\usepackage[margin=20mm]{geometry}
\usepackage{amsmath,amssymb,mathtools,fontspec,unicode-math,xurl,hyperref}
\IfFontExistsTF{DejaVu Serif}{\setmainfont{DejaVu Serif}}{\setmainfont{Latin Modern Roman}}
\setmathfont{latinmodern-math.otf}
\setlength{\parindent}{0pt}\setlength{\parskip}{6pt}\sloppy
\title{Scientific initial review of
Naked singularity censoring with anisotropic apparent horizon}
\author{}
\date{}
\begin{document}
\maketitle
Abstract: This review concerns Xinliang An’s 90-page author manuscript arXiv:2401.02003v1. All source pages were read and visually inspected. Two Major concerns are supported: an impossible signed parameterization of a nonnegative shear norm in the instability theorems, and an incorrect induced-metric calculation in the spacelikeness proof. Explicit corrections and their precise limits are given. Two local Minor concerns concern a missing background term in an expansion estimate and a false extreme Weyl-component equation. Neither the corrected horizon mechanism nor the general physical censorship conclusion is refuted by these checks. Uniformity and low-regularity existence remain evidence questions.\par
\section*{Identity, version, and assessment}
The assigned paper is Naked singularity censoring with anisotropic apparent horizon, by Xinliang An, M25-0014, associated with https://doi.org/10.4007/annals.2025.201.3.3. The actually reviewed version is https://arxiv.org/abs/2401.02003v1, dated January 2024. The manifest-bound PDF has 90 pages. Page references below use one-based PDF pages, agreeing with printed page numbers. Source-line references are one-based in the supplied 5,364-line `source.tex`. The DOI transcription amendment is metadata evidence, not a scientific finding.\par
The PDF SHA-256 is\par
> 3f27eb5f566f804046c997e903bfe6d7bca61ab8a50716e87bf805032f8e8286.\par
The TeX SHA-256 is\par
> 38c7cb1da247c255145489e0f3d55c43e1aa5ba77980be303b9d7d35086715cc.\par
All seven manifest-bound inputs were rehashed and matched. No comparison establishing equality with the journal Version of Record was performed. The findings concern this author version exclusively.\par
Assessment 1.  The overall initial judgment is Major, with two Major and two Minor supported concerns. Problem 1 contradicts the literal all-real-parameter initial-data statement, including arbitrarily small negative parameters. Problem 2 invalidates a core geometric derivation; the estimates already available suggest a verifiable repair for fixed positive barrier constants. Problems 3–4 are localized errors without demonstrated failure of the corrected core argument. Pending uniformity and regularity questions are not counted as proved defects. Independent scientific acceptance and author adoption remain pending.\par
\section*{Definitions, hypotheses, and checked scope}
Equation reference: metric. For clarity this report writes \(v=\underline u\), \(r=|u|\), \(\mu=\mu^*\in(0,1)\), and \(k=\mu/(1-\mu)\). This \(v\) is not the separately rescaled coordinate \(v=4\underline u\) introduced on source p. 20. The field equations are \[\operatorname{Ric}_{\alpha\beta}=\partial_\alpha\phi\,\partial_\beta\phi,
 \qquad \Box_g\phi=0.\] The introductory gauge on p. 2 is \[
 g=-2\Omega^2(\mathrm du\otimes\mathrm dv+\mathrm dv\otimes\mathrm du)
 +h_{AB}(\mathrm d\theta^A-d^A\mathrm dv)(\mathrm d\theta^B-d^B\mathrm dv),\] with \(e_3=\Omega^{-1}\partial_u\), \(e_4=\Omega^{-1}(\partial_v+d^A\partial_A)\), and \(g(e_3,e_4)=-2\). The shift vanishes on the initial \(H_{-1}\). Equation (2.1), p. 14, prints \(d^A\mathrm du\) instead of \(d^A\mathrm dv\). This conflicts with the introduction and its displayed frame. The analysis of Problem 2 uses [metric] (\#metric); its shift-free diagnostic is independent of this isolated coordinate discrepancy.\par
On \(v=0\), the prescribed spherical background has \[h=r^2\gamma_{\mathbb S^2},\quad
 \Omega_0^2=r^k,\quad
 \Omega^{-1}\operatorname{tr}\chi=\frac{2(1-\mu)}r,\quad
 \Omega\operatorname{tr}\underline\chi=-\frac2r,\quad
 \Omega\underline\omega=\frac{\mu}{4(1-\mu)r}.\] Theorem 1.1 combines weighted initial bounds through five angular and three outgoing derivatives, a small region \(v\le r\Omega_0^{2-\widetilde\delta}\), and an energy profile \[\int_0^v\left(|\widehat\chi|^2+
     |\Omega e_4\phi-(\Omega e_4\phi)_{\mathrm{corner}}|^2\right)\mathrm dv'
 =vf(\omega,v),\quad 0\le f\le1,\quad f\ge m>0
 \ \hbox{on a fixed ball }B_p(\epsilon).\] The parameter \(a\ge1\) is retained in later sections and may be set to one. Constants depending on fixed \(\mu,m,\epsilon\) are not uniform as \(\mu\to0,1\) or \(m,\epsilon\to0\). The elliptic unknowns are related by \[R=(va)^{1-\mu}e^{-(1-\mu)\widetilde\varphi}
   =(va)^{1-\mu}\psi^\lambda,\qquad
 \psi=e^{-\mu\widetilde\varphi}>0,\quad
 \lambda=\frac{1-\mu}{\mu}.\] The constructed MOTS is \(u=-R(\omega,v)\). Uniqueness proved in the elliptic argument uses the positive barrier and a priori bounded class. This review does not extend it to arbitrary uncontrolled graphs or to all notions of apparent horizon.\par
The entire extracted text, 14,583 lines, was read to zero remaining lines, and every source PDF page 1–90 was actually viewed. The complete TeX was machine-scanned statically; critical passages were also read verbatim. The latter does not mean every TeX line was separately read in addition to the full PDF. The coverage ledger records:\par
- pp. 1–13: all main theorem quantifiers; characteristic data; energy positivity; the degenerating \(C^0\) profile; hyperbolic/elliptic region matching and parameter dependence.\par
- pp. 14–24: metric/frame conventions; null structure and Bianchi systems; renormalized variables; derivative counts; initial norms; bootstrap smallness \(v a^{1/2}/(r\Omega^{3/2})\).\par
- pp. 25–35: transport weights, metric comparability and Sobolev use; background-subtracted expansions; Ricci-coefficient bounds; the scalar transport estimate and the nonzero background source.\par
- pp. 36–49: commutation and top angular scalar-energy pairing; cancellation of paired highest derivatives; Hodge estimates and auxiliary mass-aspect/torsion quantities.\par
- pp. 50–65: top Ricci and curvature estimates, renormalized curvature-energy pairs and closure strategy. Reciprocal references between estimates were not treated as circularity without a failed absorption argument. External existence/extension theorems were not independently reproved.\par
- pp. 66–74: MOTS equation, barrier direction, coordinate changes, natural quadratic gradient term, Moser and Schauder dependencies. The nested exponential in \(C(\kappa)\), (9.15), was retained.\par
- pp. 75–81: linearization, principal eigenfunction, maximum principle, continuation and uniqueness in the controlled class; smooth parameter dependence and the induced horizon metric.\par
- pp. 82–90: background cancellation in the instability proof; all-real amplitude quantifiers; \(C^0\) profile and intrinsic \(\dot H^{1/2}\) estimate; conformal free data and the full bibliography.\par
\section*{Problem 1: a nonnegative norm cannot carry all signed amplitudes}
\subsection*{Claim and location}
Theorem 1.3, pp. 4–5 (TeX lines 478–487), and Theorem 13.2, p. 84 (lines 5005–5032), prescribe the first component of the data as \[|\widehat\chi|(-1,v,\omega)=\sum_{i=1}^k\lambda_i f_i(v,\omega)\] for arbitrary real \(\lambda_i,\lambda'_i\), not all zero, and assert injectivity of the pair of displayed initial-data components. The proof on pp. 84–85 chooses \(f_i\ge0\), supported on disjoint discs, with \(f_i\ge m_i a^{1/2}>0\) on a smaller ball. The \(C^0\) version on p. 87 replaces this lower bound by \(m_i g(v)^{1/2}a^{1/2}>0\) for every \(v>0\).\par
\subsection*{Complete check}
Take \(k=1\), \(\lambda_1=-\varepsilon\) with any \(\varepsilon>0\), and \(\lambda'_1=0\). Then the required nonzero-parameter condition is satisfied. Choose \(v>0\) and \(\omega\) in the smaller ball supplied by the proof. All stated support, positivity and smoothness properties of \(f_1\) are retained, but \[\sum_i\lambda_i f_i(v,\omega)<0,\qquad
 |\widehat\chi|_h(v,\omega)\ge0\] for every real tensor and positive definite induced sphere metric \(h\). No such data exist. The same contradiction applies to the continuous profiles because \(g(v)>0\) for \(v>0\). Taking \(\varepsilon\) arbitrarily small rules out interpreting the obstruction as merely excessive amplitude in an upper energy bound.\par
This is an impossibility of the quantified data prescription. It is not a purported Einstein solution violating the horizon theorem; there is no admissible data tensor realizing the prescribed negative norm. Even without the proof’s nonnegative choice, an arbitrary nonzero real \(f_1\) cannot make both \(f_1\) and \(-f_1\) nonnegative everywhere. Injectivity prevents taking every \(f_i=0\).\par
\subsection*{Counterevidence and a precise repair}
Appendix A, p. 89 (lines 5230–5269), actually supplies signed free conformal data, which is the natural place for the amplitudes. In a local angular chart let \[m=\exp\!\left(\sqrt{2a}\,x
       \begin{pmatrix}1&0\\0&-1\end{pmatrix}\right),\qquad \det m=1.\] Direct differentiation, using the manuscript’s shear formula, gives \[|\widehat\chi|^2
 =\frac14\operatorname{tr}\!\bigl(m^{-1}m'm^{-1}m'\bigr)=a(x')^2,\qquad
 |\widehat\chi|=\sqrt a\,|x'|.\] The reviewer computed this exact identity independently in SymPy; no author code was executed. Thus one may set \[x'=a^{-1/2}\sum_i\lambda_i f_i\] and state the theorem for the resulting conformal metric or actual trace-free shear tensor, with the displayed norm replaced by \(\left|\sum_i\lambda_i f_i\right|\). Disjoint support still makes the squared energy a sum of \(\lambda_i^2 f_i^2\). For any nonzero shear amplitude, a corresponding smaller ball gives the required positive energy density. The scalar-field amplitudes remain signed scalar data.\par
The norm alone then loses the sign information: replacing all \(\lambda_i\) by \(-\lambda_i\) gives the same norm. Injectivity must be stated for the full free data. With \(x(0,\omega)=0\), the unique symmetric logarithm of the positive matrix \(m\) recovers \(x\), differentiation recovers \(x'\), and disjoint supports recover the coefficients. This verifies injectivity of this corrected free-data map. The characteristic constraints and the hypotheses of the evolution theorem still have to be imposed; the algebraic correction alone is not an independent existence theorem. Restricting to nonnegative coefficients would instead give a cone, not the printed all-real linear family.\par
\subsection*{Impact and judgment}
This is Major because the advertised arbitrary-dimensional instability family in Theorems 1.3 and 13.2 is a central explicit contribution, and its literal data prescription is impossible. The positive-energy mechanism and the qualitative availability of many signed conformal perturbations are substantial counterevidence against reading this as a refutation of the physical instability conclusion. Theorem 1.1 for valid initial data is not contradicted. Evidence: source pp. 84, 87, 89 and `verification/0014/evidence/independent\_checks.json`.\par
\section*{Problem 2: the spacelikeness proof uses the wrong induced metric}
\subsection*{Claim and location}
Proposition 12.4 and (12.5), p. 81, TeX lines 4851–4897, assert spacelikeness under the extra smallness condition on \(v\partial_v f\). After deriving \(R_v\ge (1-\mu)R/(2v)>0\), the proof prints \[g'_{AB}=h_{AB},\qquad g'_{Av}=2R_A,\qquad g'_{vv}=4R_v.\] These are not the pullback of the metric and frame specified on p. 2. The same paragraph also reverses the direction of a barrier bound.\par
\subsection*{Exact pullback and positivity criterion}
Equation reference: schur. Pull back [metric] (\#metric) under \((\omega,v)\mapsto(-R(\omega,v),v,\omega)\). Since \(\mathrm du=-R_A\mathrm d\theta^A-R_v\mathrm dv\), one obtains \[\begin{aligned}
 g'_{AB}&=h_{AB},\\
 g'_{Av}&=2\Omega^2 R_A-h_{AB}d^B,\\
 g'_{vv}&=4\Omega^2R_v+h_{AB}d^Ad^B.\end{aligned}\] All coefficients are evaluated on the horizon. The angular block is positive definite, so its Schur complement is \[
 4\Omega^2\left(
 R_v+d^AR_A-\Omega^2 h^{AB}R_AR_B\right).\] An independent symbolic multiplication of the full \(4\times4\) metric and the \(4\times3\) tangent matrix verified every entry and [schur] (\#schur).\par
Already at \(d=0\) and \(R_A=0\), the actual coefficient is \(4\Omega^2R_v\), whereas the printed coefficient is \(4R_v\). The paper’s interior lapse is small and is not normalized to one there. This diagnostic does not depend on the shift typo in (2.1). It refutes the displayed tensor identity, not positivity of its corrected version. No unproved realization of an arbitrary graph as a MOTS is needed for the pullback calculation.\par
For the barrier, (9.13), p. 68, gives \(\widetilde\varphi\ge\log(3/2)\), and hence \[\psi=e^{-\mu\widetilde\varphi}\le(2/3)^\mu.\] The last paragraph on p. 81 uses the reverse inequality \(\psi\ge(2/3)^\mu\). Theorem 11.6 actually places the solution strictly between the barriers. A valid positive lower bound instead comes from the upper bound \(\widetilde\varphi\le\bar\kappa(m,\epsilon)\): \(\psi\ge e^{-\mu\bar\kappa}\).\par
\subsection*{A checked repair with the required restrictions}
The corrected criterion is positivity of the bracket in [schur] (\#schur). Fix \(0<\mu<1\), \(a\), and the positive barrier constants. The preceding barrier and derivative estimates then give \[R\asymp v^{1-\mu},\quad R_v\ge cR/v,\quad
 |\nabla_\omega\log R|\le C,\quad
 h\asymp R^2\gamma_{\mathbb S^2},\quad \Omega^2\asymp R^k.\] Consequently \[\frac{\Omega^2|\nabla_hR|^2}{R_v}
 \le C\frac{\Omega^2v}{R}=O(v^{2\mu})\longrightarrow0.\] The omitted shift must also be controlled. In the paper’s frame, \(\zeta=(\eta-\underline\eta)/2\); the angular part of the frame commutator gives \[\partial_u d^A=4\Omega^2\zeta^A,\qquad d(-1,v,\omega)=0.\] The torsion estimate of Proposition 4.5 with the sphere Sobolev estimate gives \(|\zeta|_h\le Cv/r^2\), for fixed \(a\) and bounded bootstrap norms. Raising to angular coordinates supplies one further factor \(r^{-1}\). Integration from \(r=1\) to \(r=R\) yields \[|d|_{\mathrm{coord}}\le Cv\int_R^1 s^{k-3}\mathrm ds
 \le C v
 \begin{cases}
 R^{k-2},&k<2,\\
 1+\log(1/R),&k=2,\\
 1,&k>2.
 \end{cases}\] A bounded exterior contribution can be included in the constant. It follows that \[\frac{|d^AR_A|}{R_v}\le Cv|d|_{\mathrm{coord}}
 =
 \begin{cases}
 O(v^{3\mu}),&\mu<2/3,\\
 O(v^2\log(1/v)),&\mu=2/3,\\
 O(v^2),&\mu>2/3.
 \end{cases}\] Every displayed rate tends to zero for the fixed parameters. Thus the corrected Schur complement is positive for sufficiently small \(v\). This repairs the local algebraic/geometric step under the stated bounds, with data-dependent smallness and the correct lower barrier. It is not a claim of uniformity as the barrier constants degenerate.\par
\subsection*{Counterevidence, impact, and judgment}
The positive \(R_v\) bound is real evidence in favor of the intended conclusion. The lapse factor makes the angular-gradient subtraction smaller in the actual criterion. The shift has a controllable rate. Accordingly, this finding does not exhibit a timelike apparent horizon.\par
Nevertheless the printed proof establishes positivity of a different quadratic form and invokes an incorrect barrier inequality. Spacelikeness of the constructed horizon is a central geometric claim, so this invalid derivation is Major under standard 0.7.1 even though the specific local correction is feasible. The manuscript’s more general achronality statement on p. 80 also uses its external forthcoming reference [7]; that theorem was not independently read here. The repair above is for Proposition 12.4 under its additional derivative assumption, not a replacement proof of that general theorem. Evidence: source pp. 2, 68, 80–81 and `verification/0014/evidence/geometry\_checks.json`.\par
\section*{Problem 3: a background expansion term is lost at derivative zero}
\subsection*{Claim, assumptions, and contradiction}
The second estimate in Proposition 4.3, p. 30, repeated at the end of its proof on p. 31 (TeX lines 1837–1842 and 1916–1920), states \[\sum_{i\le4}\|r^i\nabla^i(\Omega\operatorname{tr}\chi)\|_{L^2(S_{u,v})}
 \lesssim \frac{va\,\mathcal O^2}{r}.\] The sum includes \(i=0\). At the prescribed \(v=0\) spherical data, with fixed \(r>0\), \[\Omega\operatorname{tr}\chi=2(1-\mu)r^{k-1},\qquad
 |S_{u,0}|=4\pi r^2,\] so \[\|\Omega\operatorname{tr}\chi\|_{L^2(S_{u,0})}
 =4\sqrt\pi(1-\mu)r^k>0.\] All higher angular derivatives of this background vanish. The asserted right-hand side is zero. The same failure persists in the one-sided limit \(v\downarrow0\) for admissible bounded-norm developments with this prescribed background. It is a test of the initial data used by the proposition, not a replacement spacetime or a parameter outside \(0<\mu<1\).\par
\subsection*{Correction and downstream check}
The proposition’s first, background-subtracted estimate is consistent: \[\sum_{i\le4}\|r^{i+1}\nabla^i
 [\Omega\operatorname{tr}\chi-\Omega^2(\Omega^{-1}\operatorname{tr}\chi)_0]\|_2
 \lesssim va\mathcal O^2.\] Together with the already stated lapse and angular-torsion bounds it gives the corrected form \[\sum_{i\le4}\|r^i\nabla^i(\Omega\operatorname{tr}\chi)\|_2
 \lesssim \Omega_0^2+\frac{va\mathcal O^2}{r}.\] The background-free estimate can also be restricted to positive angular orders, with the lapse derivatives estimated as in the proof.\par
In Proposition 5.1, p. 35, treating \(-\tfrac12(\Omega\operatorname{tr}\chi)(\Omega e_3\phi)\) as entirely \(O(v)\) loses this zeroth-order background contribution. Its weighted transport integral is nevertheless finite: the background size reduces to \(\int_r^1s^{k-1}\mathrm ds=(1-r^k)/k\), for fixed \(\mu>0\). Restoring an \(O(1)\) contribution is compatible with the required \(O(a^{1/2})\) estimate because \(a\ge1\). It does not justify the incorrectly stronger \(O(v)\) intermediate bound.\par
\subsection*{Counterevidence and impact}
Equation (4.4), p. 32, retains a \(1/r\) background term. The Raychaudhuri application on p. 83 explicitly writes \(2(1-\mu)\Omega^2/r\) back into the estimate. The decisive trapped-surface calculation uses the valid background-subtracted expansion, and the highest angular estimates remove the constant background. The supported finding is therefore a localized false estimate and its overstated intermediate source bound, not a demonstrated failure of the core bootstrap after the background is restored. Judgment: Minor. The check and corrected transport integral are analytic; the exact background norm is also recorded in the independent-check result.\par
\section*{Problem 4: the extreme Weyl equation has nonzero pure-trace matter}
\subsection*{Claim and exact location}
Equation reference: badmatter. The first line of (2.13), p. 18, TeX lines 1132–1133, presents the \(\nabla_3\alpha\) equation with matter contribution \[
 \left[\frac14(D_3\operatorname{Ric}_{44}-D_4\operatorname{Ric}_{43})
          +\frac16e_4\operatorname{Scal}\right]h_{AB}.\] Here \(\alpha_{AB}=W(e_A,e_4,e_B,e_4)\) is trace-free. All its displayed vacuum terms are trace-free as well. The scalar coefficient in [badmatter] (\#badmatter) does not vanish on all Einstein–massless-scalar solutions.\par
\subsection*{An exact solution satisfying the identity’s assumptions}
For \(t>0\), use the spatially flat scalar-field FLRW solution \[g=-\mathrm dt^2+t^{2/3}(\mathrm dx^2+\mathrm dy^2+\mathrm dz^2),\qquad
 \phi=c\log t,\quad c=\sqrt{2/3}.\] Direct calculation gives \[\operatorname{Ric}_{tt}=\frac{2}{3t^2},\quad \operatorname{Ric}_{ij}=0,\quad
 \operatorname{Scal}=-\frac{2}{3t^2},\quad \Box_g\phi=0,\quad W=0.\] Thus the exact field equations used in Section 2.2 hold. For sphere sections one can equivalently use radial spatial coordinates away from their center and take \(e_4=\partial_t+t^{-1/3}\partial_\rho\), \(e_3=\partial_t-t^{-1/3}\partial_\rho\); then \(g(e_3,e_4)=-2\). The following contractions are the same as for a unit Cartesian spatial direction at a point. Let \(H=D^2\phi\). They are \[H_{44}=-\frac{4c}{3t^2},\qquad
 H_{34}=-\frac{2c}{3t^2},\qquad
 e_3\phi=e_4\phi=\frac ct.\] Since \(D\operatorname{Ric}=D(\mathrm d\phi\otimes\mathrm d\phi)\), \[D_3\operatorname{Ric}_{44}=-\frac8{9t^3},\qquad
 D_4\operatorname{Ric}_{43}=-\frac4{3t^3},\qquad
 e_4\operatorname{Scal}=\frac4{3t^3}.\] Consequently [badmatter] (\#badmatter) equals \[\left(\frac1{9t^3}+\frac2{9t^3}\right)h_{AB}
 =\frac1{3t^3}h_{AB}\ne0.\] Every Weyl quantity in the first equation vanishes because \(W=0\), so the printed equation says \(0=h_{AB}/(3t^3)\). The trace is \(2/(3t^3)\), already nonzero at \(t=1\).\par
The reviewer independently computed the Christoffel symbols, full Ricci tensor, all 256 Weyl components, the scalar wave operator, and these covariant-derivative contractions using exact SymPy expressions. All seven grouped checks in the saved result passed. This spacetime verifies the assumptions of the universal local Einstein–scalar identity. It does not satisfy, and is not claimed to satisfy, the special naked-singularity characteristic data of Theorem 1.1. The finding is precisely the false identity in Section 2.2.\par
\subsection*{Counterevidence, correction limits, and impact}
Section 8 develops energy estimates for the renormalized \(\beta,K,\check\sigma,\underline\beta\) system rather than evolving the extreme \(\alpha,\underline\alpha\) components by this line. The static source scan and the inspected energy argument did not establish a core use of the false extreme-component equation. Judgment: Minor.\par
A valid correction requires rederiving its complete symmetric trace-free matter contribution from the divergence identity (2.12). Simply declaring the displayed pure-trace term to be projected away does not establish the missing anisotropic terms in a general scalar field. No unverified full replacement equation is proposed. The unused incorrect identity can instead be removed without changing the inspected renormalized energy argument. Evidence: source p. 18, Section 8, the independent script and its result.\par
\section*{Unresolved evidence questions and excluded concerns}
\subsection*{Uniform smallness in the angular concentration}
The main statements write \(\delta=\delta(B)\). The barrier construction and region match in (1.15) and (9.16), pp. 12 and 68, visibly involve \(m,\epsilon\). For example put \(s=m\epsilon^{\tau+2}\), \(a=1\), and \(c_0=\mu\widetilde\delta/2\in(0,1)\). A required smallness inequality of the form \(s\ge(vs)^{c_0}\) is equivalent to \[v\le s^{(1-c_0)/c_0}.\] There is no fixed positive \(v\) satisfying this as \(s\downarrow0\). The a priori \(C^2\) barrier constants also deteriorate with \(m,\epsilon\). The manuscript should specify whether the horizon interval is shrunk after these data are fixed and whether that shrinkage is included in the notation \(\delta(B)\).\par
This proves a limitation of the displayed proof’s uniformity. It does not prove nonexistence of MOTSs for concentrated data beyond that sufficient range. Classification: Needs Evidence, not an additional proved Major issue.\par
\subsection*{Weighted norms, the auxiliary initial norm, and continuous data}
Page 23 defines \(\mathcal I^{(0)}\) using unweighted outgoing derivatives through order five and angular derivatives through order six. Page 24 permits implicit constants to depend on initial-data quantities, while the theorem’s displayed controlled norm has only three weighted outgoing derivatives and five angular derivatives. The continuous example on pp. 85–88 uses \[h(v)=[\log\log(1/v)]^{-1/4},\qquad
 h'(v)=\frac{1}{4v\log(1/v)[\log\log(1/v)]^{5/4}}.\] Thus \(h'(v)\to\infty\), but \(v^j h^{(j)}(v)\to0\) for \(j=1,2,3\). The latter three limits were symbolically checked. This satisfies the intended pointwise weighted-derivative behavior, not the unweighted \(\mathcal I^{(0)}\) bound.\par
There is also a notational ambiguity: some displayed theorem expressions place \(v^j\) outside an \(L^\infty_v\) norm. If that is taken literally as a multiplier after an infinite unweighted supremum, the continuous example fails it. The function space should explicitly put the weight inside the supremum.\par
Even in a smooth one-sided family, low derivative control does not by itself control five outgoing derivatives: a profile \(C+b(1+N^{-3}\sin Nv)\), with fixed small \(b>0\), has uniformly bounded derivatives through order three on a fixed interval, whereas its fifth derivative has size \(bN^2\). It can retain a uniformly positive, bounded energy density after normalization. This elementary family shows that a bound for \(\mathcal I^{(0)}\) cannot simply be inferred from the displayed low-order norm. It is not a construction of a failed Einstein evolution.\par
To complete the continuous-data assertion, one needs an existence and limiting argument whose lifespan and necessary estimates are uniform in smooth approximations of these profiles, or a direct theorem in the specified weighted function space. The manuscript’s many estimates use angular rather than high outgoing derivatives, which is genuine evidence that such a repair may be possible. The reference to standard local existence and extension on p. 65 alone does not record that uniform argument. This remains Needs Evidence; no failure of the claimed low-regularity solution has been proved here.\par
\subsection*{Checks that do not support additional findings}
Proposition 13.4 estimates the intrinsic Slobodeckij seminorm on \([0,\delta]\). Its near/far diagonal decomposition gives the stated smallness; the logarithmic substitutions and integration bounds were checked. The zero extension to \(v<0\) would instead have cross energy comparable to \(\int_0^\delta h(v)^2\,\mathrm dv/v=\infty\). Because the proposition explicitly uses the one-sided interval, this does not refute it. The intrinsic result must not be advertised as a global critical norm bound across the junction without further work.\par
The linearized principal term is \[\Delta_\gamma-\mu-\frac{1-\mu}{2}f\psi^{-1/\mu}.\] For fixed \(\mu>0\), its negative zeroth-order contribution remains coercive even where \(f=0\). Thus the p. 86 informal global notation \(f\psi^{-1/\mu}\asymp1\) is too strong outside the support, but a positive lower bound everywhere is not needed in this argument. The positive barrier context must be retained when lemmas say nonnegative \(\psi\), because negative powers are used.\par
The nested bound (9.15) is \(C(\kappa)=(Ce^\kappa)^{C e^{4\Lambda\kappa e^\kappa}}\). For the log-log profile, fixed angular-data constants add to \(\kappa\); they cannot be discarded in exact equalities. Nevertheless \(4\Lambda\kappa e^\kappa=o(\log\log(1/v))\) as \(v\downarrow0\) for that profile, so the qualitative domination by arbitrarily small powers of \(v^{-1}\) survives after shrinking the interval. No extra false-conclusion finding is supported.\par
Isolated coordinate, sign and numerical-normalization inconsistencies were not counted without a demonstrated substantive dependency. The manuscript’s external reference [7] is listed as in preparation. Its absence from this review is a dependency limitation, not a missing-citation finding. Neither global future null infinity nor a global formulation of cosmic censorship was independently established by this local review.\par
\section*{Dependency, outcome table, and remaining work}
The main proof architecture is weighted hyperbolic control, then barriers and elliptic continuation, then a family of MOTSs, followed by geometric and instability conclusions. Problem 1 affects the literal family of prescribed norm data. Problem 2 affects the geometric proof of spacelikeness under the stated extra derivative condition. Problem 3 is a background bookkeeping error with a checked replacement sufficient for the inspected use. Problem 4 is an unused extreme curvature equation in the inspected dependency chain. None of these checks establishes a counterexample to the corrected existence of anisotropic trapped surfaces.\par
Item | Judgment | Repair verification | Remaining work\par
P1 | Major | Signed conformal data and absolute norm; full-data injectivity checked | Restate theorem and maintain constraint/evolution hypotheses\par
P2 | Major | Exact Schur complement; fixed-parameter positivity from torsion and barrier bounds | Insert corrected proof; general achronality dependency\par
P3 | Minor | Restore \(\Omega_0^2\); finite background transport integral | Propagate corrected intermediate bounds\par
P4 | Minor | Counterexample exact; no complete replacement claimed | Rederive STF equation or remove unused line\par
Uniformity | Needs Evidence | Sufficient smallness dependence derived | Specify dependence on concentration\par
Regularity | Needs Evidence | Weighted/unweighted distinction checked | Uniform limiting existence argument\par
No repair has been submitted to or adopted by the author. No independent later-version or erratum match was performed. No empirical experiment is needed to verify the algebraic contradictions; the outstanding work is analytic, version-comparative, and editorial. A journal-final defect cannot be inferred until the exact corresponding Version-of-Record passages are checked.\par
\section*{Actual computation and static inspection}
Two reviewer-authored scripts, run with `[local archive path omitted] `-X utf8 -B`, are preserved under `verification/0014/evidence/`. `independent\_checks.py` ran as PID 28408, exit code 0, using SymPy 1.13.3. Its result contains:\par
- exact Ricci/scalar-wave equations and 256 zero Weyl components;\par
- nonzero printed \(\alpha\)-matter coefficient \(1/(3t^3)\);\par
- Appendix A shear norm squared \(a(x')^2\);\par
- the three vanishing weighted profile-derivative limits and the divergent unweighted first derivative.\par
`geometry\_checks.py` ran as PID 26748, exit code 0. It verified the pullback blocks, Schur complement, the \(2\mu\) and \(3\mu\) powers, and the direction of the constant barrier bound. The computations are exact algebra and limits, not numerical Einstein evolution, empirical reproduction, or independent PDE existence proofs.\par
The static scan read all 5,364 source lines and recorded targeted identity/dependency hits. No executable author code was run. Its inline bibliography and diagrams were inspected through the original PDF. Directly read critical source ranges include 317–351, 979–1078, 1121–1161, 1439–1546, 1833–1922, 4847–4905, 5001–5052, and 5219–5272.\par
\section*{Supporting checks}
The supplied PDF was extracted with the installed Poppler tools. The extraction completed as PID 17144, with 90 pages. The renderer completed as PID 25232, with 90 PNGs; a Symbol-font display warning was recorded. Every rendered source page was actually viewed. The source text ledger ends at line 14,582 using zero-based indexing, with zero remaining lines.\par
No `pdflatex`, `xelatex`, `lualatex`, or `tectonic` was found on PATH or the recorded standard installation paths. No package or compiler was installed. Therefore `reports/Pdf/0014.pdf` does not exist as a compiled report deliverable, and no report-page visual QA is claimed. A static report-environment check does not replace compilation.\par
The source manifest SHA-256 is\par
> 75bfd0081a43686f6b6ebf79fcc143e14debfeddb39c6a3b83fe503f1d8fb553.\par
The scientific review, mathematical derivations, reviewer code, static inspection and page inspection were performed by this assistant. No scientific subagent or human scientific approval is claimed.\par
This assistant did not personally inspect the dispatcher’s browser tab. These are separately attributed observations; no equivalence or backend route verification is invented. No independent evidence of a non-6-Pro route was observed, and no model or mode switch was made.\par
9 X. An, Naked singularity censoring with anisotropic apparent horizon, arXiv:2401.02003v1, January 2024, 90-page supplied author manuscript. https://arxiv.org/abs/2401.02003v1. The exact source version, not an independently matched journal version, is the primary scientific reference consulted. Annals of Mathematics, article identity record, DOI 10.4007/annals.2025.201.3.3, https://annals.math.princeton.edu/2025/201-3/p03, supplied acquisition snapshot of 25 October 2026. Used for identity, not for an unperformed text equivalence claim.\par
\section*{Compilation continuation (26 October 2026)}
Scientific conclusions, pending questions, source-version limits and the absence of independent acceptance remain unchanged.\par
\end{document}
